\documentclass[11pt]{article}
\usepackage[margin=1.15in]{geometry}
\usepackage{amsmath,amssymb,amsthm}
\newtheorem{proposition}{Proposition}

\title{How Long Does Tea Stay Warm?\\\large A Toy Model of Cooling in a Mug}
\author{Alice Moss \and Bob Okafor}
\date{November 2026}

\begin{document}
\maketitle

\begin{abstract}
We model a mug of tea as a single well-mixed body that loses heat through its surface and its walls.
The model has two parameters, both easy to measure in a kitchen, and predicts the time until the tea
drops below a drinkable temperature to within a minute in our tests.
\end{abstract}

\section{Introduction}
Everyone who has forgotten a cup of tea on a desk knows the answer is ``not long enough''.
Here we ask how long, exactly, and which choices matter most: the shape of the mug, a lid, or milk.

\section{Model}
Let $T(t)$ be the temperature of the tea and $T_a$ the temperature of the room. Newton's law of cooling gives
\begin{equation}
  \frac{dT}{dt} = -k\,(T - T_a), \qquad T(0) = T_0,
  \label{eq:newton}
\end{equation}
where $k = hA/(mc)$ depends on the exposed area $A$, the mass $m$ and the heat capacity $c$.
The solution is $T(t) = T_a + (T_0 - T_a)\,e^{-kt}$.

\begin{proposition}
The tea stays above a temperature $T_d$ for a time
\[
  t_d = \frac{1}{k}\,\ln\frac{T_0 - T_a}{T_d - T_a}.
\]
\end{proposition}

\section{Results}
With $T_0 = 90^\circ$C, $T_a = 21^\circ$C and $T_d = 55^\circ$C, a lid halves $A$ and so doubles $t_d$.
Adding milk lowers $T_0$ but also raises $m$; in our tests the two effects nearly cancel.

\begin{table}[ht]
  \centering
  \begin{tabular}{lccc}
    Mug & $A$ (cm$^2$) & $k$ (1/min) & $t_d$ (min) \\
    \hline
    Ceramic, open & 50 & 0.031 & 22 \\
    Ceramic, lid & 25 & 0.017 & 41 \\
    Glass, open & 50 & 0.036 & 19 \\
    Steel, double wall & 50 & 0.024 & 29 \\
  \end{tabular}
  \caption{Measured cooling rates and drinkable times for four mugs.}
  \label{tab:mugs}
\end{table}

\section{Discussion}
The single-body model ignores the temperature gradient inside the tea. Near the surface the tea cools
faster than in the middle, and a sip taken from the top is cooler than the average $T$ predicts.
Stirring removes most of the difference, which is one reason the model fits stirred tea best.

Evaporation is the largest loss we leave out. At $90^\circ$C it carries away more heat than
radiation and conduction through the walls together, and it falls off steeply as the tea cools.
A lid stops most of it, which explains why the lid does better than halving $A$ alone would suggest.

The walls matter less than people think. A double-walled steel mug lowers $k$ by about a fifth,
while a lid lowers it by almost half (Table~\ref{tab:mugs}).

\section{Limitations}
We measured one mug of each kind, in one kitchen, in November. The room was at $21^\circ$C and
still; a draught raises $h$ and so $k$. Our thermometer read to $0.5^\circ$C, which limits $t_d$
to about a minute either way.

\section{Conclusion}
A lid is the cheapest way to keep tea warm: it nearly doubles the time the tea stays drinkable.
Everything else, from the mug's material to a splash of milk, moves $t_d$ by a few minutes at most.

\appendix
\section{Fitting $k$}
We fit $k$ by least squares to $\ln(T - T_a)$ against $t$, which is linear under \eqref{eq:newton}:
\[
  \ln\bigl(T(t) - T_a\bigr) = \ln(T_0 - T_a) - k\,t .
\]
Each fit used readings every 30 seconds for 40 minutes. The residuals were below $0.4^\circ$C
except in the first two minutes, while the mug itself was still warming up.

\section{Measurements}
For each mug we poured 250 ml of water at $90^\circ$C, waited for the reading to settle, and logged
the temperature every 30 seconds. The room temperature was logged at the same times.
We repeated each measurement three times and report the mean $k$; the spread was below 4\%.

\end{document}
